% Generated draft: requires mathematical and source review.
\documentclass[11pt]{article}
\usepackage[margin=25mm]{geometry}
\usepackage{fontspec}
\setmainfont{Latin Modern Roman}
\usepackage{amsmath,amssymb,mathtools,unicode-math}
\setmathfont{Latin Modern Math}
\usepackage{enumitem,needspace,xurl,hyperref}
\setlength{\parindent}{0pt}
\setlength{\parskip}{5pt}
\setlength{\emergencystretch}{3em}
\newcommand{\N}{\mathbb{N}}
\newcommand{\Z}{\mathbb{Z}}
\newcommand{\Q}{\mathbb{Q}}
\newcommand{\R}{\mathbb{R}}
\newcommand{\C}{\mathbb{C}}
\newcommand{\F}{\mathbb{F}}
\newcommand{\A}{\mathbb{A}}
\newcommand{\PP}{\mathbb{P}}
\newcommand{\E}{\mathbb{E}}
\newcommand{\eps}{\varepsilon}
\newcommand{\dd}{\,\mathrm{d}}
\begin{document}
\section*{UnsolvedMath ID: 30000571}
\textbf{Source problem code:} catalogue entry
\subsection{Definitions and mathematical statement}
Let $T$ be a hereditarily hypercyclic operator in the space and scalar-field setting intended by the source. For a bounded linear operator on a Banach space $X$, use the standard notation
\[
 HC(T)=\{x\in X:\overline{\{T^mx:m\in\N_0\}}=X\}.
\]
One common hereditary-hypercyclicity convention says there exists an increasing sequence $(m_j)$ such that, for every subsequence $(m_{j_k})$, the family $(T^{m_{j_k}})$ has a dense orbit. The source does not state its convention, and a condition on every subsequence of all powers would have a different quantifier pattern.

The retained first question asks whether $HC(T)$ contains the range of an injective bounded operator $S$. Its literal wording needs correction or a convention: every linear range contains zero, whereas zero is not hypercyclic on a nonzero space. A possible intended formulation is
\[
 \operatorname{Im}(S)\setminus\{0\}\subseteq HC(T),
\]
or a convention that writes $HC(T)$ for the hypercyclic vectors together with zero. The record does not authorize choosing between them. It also omits the domain of $S$ and whether $S$ should act on $X$ itself or another infinite-dimensional Banach space. No closed-range condition is added.

The second question asks whether every linearly independent $n$-tuple in the same range is hypercyclic for the $n$-fold direct sum. The target property is
\[
 x_1,\ldots,x_n\in\operatorname{Im}(S),\quad
 x_1,\ldots,x_n\text{ linearly independent}
 \quad\Longrightarrow\quad
 \overline{\{(T^mx_1,\ldots,T^mx_n):m\in\N_0\}}=X^n.
\]
The single common iterate $m$ matters; individual hypercyclicity of each coordinate is not this joint property. The source does not say whether it seeks one $S$ with this property for every positive integer $n$, a chosen $S$ for each $n$, or an automatic conclusion for every witness to the first question. These possibilities are not silently identified.

Thus both questions are retained, with the zero-vector issue and missing domain and quantifiers explicit. No existence of an operator range or joint-hypercyclicity conclusion is asserted.
\subsection{Short English statement}
The source asks for a bounded injective operator range of hypercyclic vectors and for simultaneous hypercyclicity of independent tuples from it. The range necessarily contains zero, so the usual nonzero-vector convention needs explicit clarification, as do the domain of $S$ and the quantifiers in the tuple question.
\subsection{Sources}
\begin{itemize}
\item Local UnsolvedMath registry, permanent ID 30000571, code catalogue entry, ``Linear Subspaces of Hereditarily Hypercyclic Vectors.'' Proposed year 2006; proposer unspecified.
\item \emph{Mini-Workshop: Hypercyclicity and Linear Chaos} (2007), Oberwolfach Reports, pp.\ 2227--2276, DOI 10.4171/owr/2006/37; \url{https://doi.org/10.4171/owr/2006/37}. The displayed report year 2007 differs from the DOI and proposed year 2006. This is the sole supplied literature reference.
\item Statement status corrected\_verified; the supplied note says operator, hypercyclic-set, tuple, image, and direct-sum notation was enclosed in balanced delimiters, retaining both questions.
\item Literature triage checked 2026-08-21: open. It generically calls the item an unresolved problem/conjecture/direction and reports no authoritative resolution in the supplied targeted check. Created 2026-08-25; updated 2026-09-16. Exact conventions and current claims require independent verification.
\item Registry curation by Ulam AI / UnsolvedMath Contributors, CC BY 4.0; the underlying report retains its supplied provenance. Unreviewed statement-only draft.
\end{itemize}
\end{document}
